\begin{helper}[Restoration induction for Algorithm (A)]
Let \(\mathcal R\) be a decreasing
\noderef{ExtensionBaseRestorationTrace} starting at a proper coloring
\(\sigma_n\) which leaves exactly \(e_0\) uncolored. Fix a distinguished
support vertex \(v_*\) in the final tree and a distinguished, exempt slot
\(q\) whose exact chosen vertex satisfies
\(\operatorname{chosenVertex}(q.\mathrm{succ})=v_*\). Assume the order
and initial-segment conclusions for every restoration slot: whenever \(t\)
has exact support \(v_t\), the distinguished support \(v_*\) precedes
\(v_t\) in the final tree order, and \(T_j\) is an ordered initial segment
of \(T_t\). Assume initial strong interchangeability of the exact
gamma/delta pair at every active slot \(t\ne q\). Assume also the
active-pair overlap law (O): for active \(s<t\), their exact stored pairs
overlap if and only if their supports agree and no member of their full
\noderef{ExtensionBaseSupportClass} lies strictly between them, and in that
case their intersection is both \(\{\gamma_t\}\) and
\(\{\delta_s\}\). Suppose finally that, for every trace step \(a<r\),
the trace's stored delta and gamma colors satisfy
\[
 \delta_{t_a}\in\overline{\sigma^a}(v_{t_a})
 \quad\hbox{and}\quad
 \gamma_{t_a}\in\sigma^a(v_{t_a}).
\]
Thus the stored delta color is missing at the trace support under the
pre-step trace coloring, while the stored gamma color is present there
under that same coloring.
Then \(\mathcal R\) satisfies the
\noderef{ExtensionBaseRestorationInvariant}.
\end{helper}
\begin{proof}
We induct on the trace position. At position zero, the trace's initial
coloring is \(\sigma_n\), so properness and the exact-one-uncolored property
are the two initial assumptions. At any later position they follow from the
preceding switch, as explained below.

Fix a step \(a\), write \(t=t_a\), \(v=v_t\),
\(\gamma=\gamma_t\), and \(\delta=\delta_t\). The enumeration clause of
\noderef{ExtensionBaseRestorationTrace} gives \(t\in L^*\), and hence, by
\noderef{ExtensionBaseRestorationIndices} and
\noderef{ExtensionBaseActiveIndices}, \(t\) is an active PE slot. Inspect the
certificate at \(t+1\) in \noderef{TashkinovSeries}. Since this is a positive
slot of exact PE kind, only the PE constructor of
\noderef{TashkinovSeriesStep} is possible. The trace's exact chosen-vertex
field identifies \(v\) with the maximum defective-choice vertex in that
constructor. The vertex-defective field, expanded through
\noderef{DefectiveVertex}, gives
\[
                         v\in V(T_t).                \tag{1}
\]
Membership \(t\in L^*\) also gives \(t\ne q\): if \(t=q\), the
distinguished-slot exact chosen-vertex hypothesis and the trace field would
put \(t\) in the removed distinguished-support class. For the first step,
initial strong
interchangeability therefore supplies the pair at \(t\). For a later step,
the preceding induction state supplies it, since strict decrease ensures
that \(t\) did not occur at an earlier position. The assumed color state
supplies that \(\delta\) is missing and \(\gamma\) present at \(v\).
Consequently \noderef{ExtensionBaseAlgorithmASingleStep} gives a nontrivial
path component, proves that it meets \(V(T_t)\) exactly in \(v\), and proves
properness and exact-one-uncoloredness after the switch. The trace transition
identifies that switched coloring with the next trace coloring. The required
distinguished-support order and \(T_j\)-initial-segment clauses are the two
standing assumptions. Together with (1), this proves every non-preservation
clause of the invariant at step \(a\).

It remains to preserve all unprocessed active pairs at slots \(s\ne q\).
Let \(s\) be such a slot, with exact support \(v_s\) and exact pair
\(\{\gamma_s,\delta_s\}\). The induction hypothesis supplies its strong
interchangeability before the current switch. If \(v_s\ne v\), then
\(s\ne t\); orient the two distinct finite ordinals increasingly and apply
the overlap law (O). Its support-equality condition is false, so the two
stored pairs are disjoint.

Suppose instead that \(v_s=v\). Since \(t\in L^*\), its support differs from
the distinguished support. Thus active membership puts \(s\) in the same full
support class and hence in \(L^*\). The trace enumerates \(L^*\), so \(s\)
occurs somewhere. It has not occurred at or before \(a\), by the unprocessed
hypothesis. Strict decrease of the trace then gives \(s<t\). We have
therefore proved in the two cases exactly the alternative required by
\noderef{StrongInterchangeabilityKempePreserved}: different supports with
disjoint color pairs, or equal supports with \(s<t\). That node, applied
with the current nontrivial path and the two strong-interchangeability
certificates, proves that the pair at \(s\) remains strongly interchangeable
after the switch. This completes the induction step.

The induction supplies properness and exact-one-uncoloredness at all
positions, every step clause, and preservation of every active slot \(s\ne q\)
absent through the current step. At the final position, the same
induction statement gives the final absent-slot clause of
\noderef{ExtensionBaseRestorationInvariant}. If the trace has no steps,
there is no step clause to prove, its sole coloring is \(\sigma_n\), and the
final absent-slot clause is exactly the initial strong-interchangeability
assumption. Hence the invariant also holds in the empty-index case.
\end{proof}
